\documentclass[10pt,a4paper]{amsart}
\usepackage[margin=19mm]{geometry}
\usepackage{amsmath,amssymb,mathtools}
\usepackage[scr=rsfs,cal=euler]{mathalfa}
\usepackage{xfrac}
\usepackage[unicode,hidelinks,bookmarks=false]{hyperref}
\newcommand{\ie}{i.e.\ }
\newcommand{\cf}{cf.\ }
\newcommand{\resp}{respectively\ }
\newcommand{\Subsection}{\subsection}
\providecommand{\nobreakdash}{\nobreak\hbox{-}}
\setlength{\emergencystretch}{2em}
\multlinegap0pt
\usepackage{bm}
\usepackage{bbm}
\usepackage{dsfont}
\usepackage{xspace}
\usepackage{cancel}
\usepackage{mathtools}
\usepackage{amsthm}
\makeatletter
\@ifundefined{prop}{\newtheorem{prop}{Proposition}}{}
\makeatother
\title[Each generic polytope in $\mathbb{R}^3$ has a point with ten]{Each generic polytope in $\mathbb{R}^3$ has a point with ten normals to the boundary}
\author{Ivan Nasonov, Gaiane Panina}
\date{Source: arXiv:2411.12745v4 (CC BY 4.0)}
\begin{document}
\maketitle

\noindent\emph{Source and licence.} Each generic polytope in $\mathbb{R}^3$ has a point with ten normals to the boundary. Version arXiv:2411.12745v4 (CC BY 4.0), distributed under the Creative Commons Attribution 4.0 International licence, as recorded in the arXiv licence metadata for this exact version and shown on the abstract page https://arxiv.org/abs/2411.12745v4. Full rights notice: licenses/arxiv-2411.12745-CC-BY-4.0.txt.\par\smallskip

\noindent\textit{Editorial scope.} Counted unit: the Prop whose source label is \texttt{propNice} in the article (source0.tex lines 355--359, source label \texttt{propNice}), delivered together with the article's own complete original proof of it (source0.tex lines 361--411, one \texttt{proof} environment of 51 lines). No mathematics is written, repaired or re-derived: statement and proof are the article's own text line for line. Required context retained uncounted: none. Every definition, display and label of those lines is retained unchanged. Omitted from this package and not redistributed: 1--354, 360--360, 412--794. Nothing inside a retained excerpt is omitted or altered: every retained body is delivered exactly as the article's lines are written, so no omission receipt of any kind accompanies this package. Citations: the retained excerpts carry no citation command at all, so no bibliography entry of the article is transcribed and no reference is redistributed. Reference adaptation: none was needed and none was made; every reference inside the delivered unit resolves inside it, so no unresolved reference or citation is produced. Printed slips kept literal: no printed word, symbol, hypothesis or number of the counted statement or of its proof was altered. Layout and class: the article's own class is \texttt{amsart} with packages including graphicx, fontenc, inputenc, babel, amsmath, amssymb, amscd, amsfonts, amsthm, xy, xfrac, color; the wrapper is the lane's pinned \texttt{amsart} editorial wrapper at 10pt on A4, which restates the theorem-like environments the retained text needs and adds the article's own macro definitions that the retained text needs (none), with extra wrapper packages (bm, bbm, dsfont, xspace, cancel, mathtools). The delivered numbering is the unit's own. Assets: no retained span contains a figure environment, a graphics-inclusion command, a PDF-page inclusion, a table or a TikZ picture; no image file is copied into this package, the package declares no asset dependency and no image file is redistributed with it. Licence: version arXiv:2411.12745v4 of this article is distributed under the Creative Commons Attribution 4.0 International licence, as recorded in the arXiv licence metadata for this exact version and shown on the abstract page https://arxiv.org/abs/2411.12745v4. A full rights notice is delivered at licenses/arxiv-2411.12745-CC-BY-4.0.txt. Characters: every retained excerpt is pure ASCII, so the delivered engine, XeLaTeX with its default Latin Modern text font, reports no missing character.
\begin{prop} \label{propNice}

If a generic polytope $\mathbf{P}\subset \mathbb{R}^3$ has no point in its interior with 10  emanating normals, then all the vertices of  $\mathbf{P}$ are skew.

\end{prop}

\begin{proof}

Assume the contrary: $V$ is a nice vertex of $\mathbf{P}$, and $\mathbf{P}$ has no point in its interior with 10  emanating normals. Let $Y$ be a point inside $P = Spher(V)$ with three short maxima $A_1, A_2,A_3$. Each of the maxima is surrounded on $\partial P$ by two short minima, so there are at least three short minima of $SQD^{sph}_Y$.  Consider two cases. 

\begin{enumerate}
    \item The number of the short minima (that are neighbors of $A_i$) is exactly $3$. Then all the minima are short, \iffalse and moreover,
 $|Yp|<\pi/2 \ \ \forall p\in \partial P$,\fi  
 $Y\in Core(Spher(V))$, so $Y\in \mathcal{AR}(V)$.

     Therefore,
for every point on the ray $VY$, the vertex $V$ is a  maximum. Take a point $y$ lying
on the ray very close to $V$ and let it travel along the ray. We may assume that $y$ never crosses two sheets of the bifurcation set at a time. If this is not the case, perturb $y$ a little bit. 

 Initially $SQD_y$ 
has  three minima $m_1$, $m_2$, and $m_3$, lying on the faces incident
to $V$ , three saddles $s_1$, $s_2, s_3$, lying on edges that correspond to $A_1, A_2, A_3$, and one maximum $V$.  Since $V$ is not
 the global maximum, there is at least one other maximum. Altogether there are at least 8 critical points of $SQD_y$, or, equivalently,
8 normals emanating from $y$.

The point $y$ keeps moving until the \textit{first event},  which is either a bifurcation of $SQD_y$,  or  the point $y$ leaves the
interior of $\mathbf{P}$. We have the following case analysis:




(a) The first event is a birth of two new critical points. This gives us 10 normals.

(b)
The point $y$ leaves the polytope $\mathbf{P}$. This means that $y$ hits one
of the faces which is not incident to the vertex $V$. Right before
$y$ leaves the polytope $\mathbf{P}$, there is the  fourth minimum on this face, and we have at least 10 critical points.

(c) One of the saddles $s_i$ dies. This happens in two cases:
(i) $s_i$ meets a  maximum of $SQD_y$. The maximum cannot be $V$ since  $s_i$ slides away from $V$.
It is not any other maximum $V'$. Indeed, recall
that $s_i$ is the orthogonal projection of $y$ to the edge. If $s_i$
meets $V'$ before $y$ leaves $\mathbf{P}$, then $|Vy|>|VV'|$. Then $V'$ can't be maximum of $SQD_y$ when $y=V$.
(ii) $s_i$ meets a minimum. It can be none of $m_1$, $m_2$, and $m_3$, so
right before the bifurcation, the function $SQD_y$ has $4$ minima and $2$ maxima, which proves the claim.

\item The number of short minima that are neighbors of $A_i$ is at least $4$. 

 We repeat the above construction. Initially we have at least four minima of $SQD_y$, three saddles, and  one  maximum.  The bifurcations are analyzed in the same way. If there is no bifurcation until $y$ leaves $\mathbf{P}$, we arrive at a fifth minimum just before leaving the polytope, which gives $10$ normals.



    
\end{enumerate}

    
\end{proof}

\end{document}
